\chapter*{\textbf{ANEXO A – TÍTULO}} %\textcolor{red}{(OPCIONAL)}}}
\addcontentsline{toc}{chapter}{ANEXO}


Nos anexos são colocados textos ou documentos não desenvolvidos pelo autor, mas que podem auxiliar na fundamentação e/ou comprovação do que foi argumentado no trabalho. Os anexos são identificados por letras maiúsculas consecutivas, travessão e seus respectivos títulos. Podem ser utilizadas letras maiúsculas dobradas quando o número de anexos esgotar as 26 letras do alfabeto.\\

\noindent
\textbf{Exemplos:}\\

\textbf{ANEXO A} – Representação gráfica da contagem de células inflamatórias – Grupo II

\textbf{ANEXO B} – Representação molecular

….

\textbf{ANEXO AA} – Representação gráfica da contagem de células inflamatórias – Grupo XXVII

